\makeatletter
\def\input@path{{./apqdesign}}
\makeatother

\DocumentMetadata{
	pdfstandard=a-2b,
	lang=en,
	pdfversion=1.7,% 2.0 is possible as well, but some validators don't support that yet
}

%Ab Hier kommen die Einstellungen für das NEUE APQ-TUD_Design!!!
\documentclass[
class=book,
accentcolor=1b,
custommargins=geometry,
fontsize=11pt,
thesis={type=master, reviewer-on-uppertitleback=true}, %options: bachelor, master, proposal, dr and drfinal
ruledheaders=all,
headline=false,
instbox=false,
marginpar=false,
title=small,
ignore-missing-data=true,
twoside=true, %for bachelor, master, dr, and drfinal: true ; for proposal: false
pdfx=false, % Disabled anyway when using \DocumentMetadata
IMRAD = false,
logofile=apqdesign/tuda_logo.pdf
]{apqpub}


% % % % % % % %
%Pakete
\usepackage[main = british]{babel} %main legt die Hauptsprache fest (zb von oder by)
\usepackage{csquotes}
\usepackage{cancel}
% package ulem ändert das Design des Literaturverzeichnisses bei Osa_Mod. Für durchstreichen kann cancel verwendet werden.
\usepackage{float}
%\usepackage{subfig}
\usepackage{amsmath,amsthm,amssymb}
\usepackage{mathtools}
\usepackage{enumerate}
\usepackage{multirow}
\usepackage{isotope}
\usepackage{subcaption}
\usepackage{graphicx}
\ifLuaTeX
	%no microtyping option for XeTeX
	\usepackage[%
	activate={true,nocompatibility},%
	tracking=true,%
	factor=1100%
	]{microtype}
	\SetTracking{encoding={*}, shape=sc}{40}
\fi
\ifPDFTeX
	\SetTracking{encoding={*}, shape=sc}{40}
    \usepackage[utf8]{inputenc}
\fi
\usepackage{nicefrac}
\usepackage{upgreek}
\usepackage[section]{placeins}
\usepackage{braket}
\usepackage{chngcntr}
\usepackage{mathrsfs}
\usepackage{adjustbox}

\usepackage[abbreviations,automake]{glossaries-extra}
\makeglossaries
\input{Content/Glossary}

\counterwithout{figure}{section}
\usepackage{siunitx}
\sisetup{
detect-mode,
  mode = text,
  locale = US ,
  per-mode = fraction,
%  per-mode = symbol,
  separate-uncertainty = false,
}
\DeclareSIUnit{\decibelmilliwatt}{dBm}
\usepackage{geometry}
\geometry{
    reset,
    a4paper,
    top=20mm,
    bottom=25mm,
    inner=31mm,%7mm Bindekorrektur sind jetzt
    outer=24mm,%der neue Standard im APQ-Design!
    includefoot,
    includehead,
    footskip=10mm,
    nomarginpar
}

\usepackage[backend=biber, sorting = none, maxnames=99, minbibnames = 2, giveninits=true, doi = false, isbn = false, date = year, url = false, abbreviate = false, bibstyle = numeric-apq, citestyle=numeric-comp]{biblatex}
\addbibresource{mybib.bib}

\usepackage{hyperref} % Immer als letztes laden
\def\UrlFont{\normalfont}
\hypersetup{colorlinks=true,linkcolor=black,urlcolor=TUDa-1b,citecolor=TUDa-1b}
\urlstyle{same}

% % % % % % % %
%Weitere Änderungen
\DeclareOldFontCommand{\bf}{\normalfont\bfseries}{\mathbf}
\RedeclareSectionCommand[afterskip=\baselineskip]{chapter}
\newcommand\todo[1]{\textcolor{red}{#1}}
\parindent 0pt
\addtokomafont{caption}{\rmfamily\small}
\addtokomafont{captionlabel}{\accentfont}

\addTitleBoxLogo*{\makebox[\linewidth][l]{\includegraphics[width=55mm]{graphics/titel_apq-logo.pdf}}}

%Möglichkeit einen generischen Pfad zu Abbildungen anzugeben.
\graphicspath{{C:/Users/Max_Mustermann/Uni/Master/Masterarbeit/Results/}}

\lowertitleback{
\captionof{figure}{$3\times3$ grid of multi-tone Top Hat profiles created with an \gls{aod}.}
\label{fig:Gesamt}
}



\begin{document}
    %\renewcommand*{\figurename}{Abb.}
    \renewcommand*{\figurename}{Fig.} %Für englische Dokumente auskommentieren
    \renewcommand*{\tablename}{Tab.}
    %\titleimage{\includegraphics[width=\textwidth, height=1.198\textwidth]{graphics/title_page.pdf}} %! Watch out! The image/pdf will be scaled to 1:1.198. Only use images with the proper aspect ratio to avoid distorted images.
    % Titelbild aus BeamShaperPlots (Reiter "Titelbild"). Die Stern-Variante skaliert
    % das Bild auf das Fenster zwischen den beiden Strichen und beschneidet es mittig
    % -> kein weisser Rand, auch wenn sich der Titelblock aendert.
    \titleimage*{\includegraphics{graphics/Titelbild.pdf}}
    \Metadata{
       title= Titel der Thesis, %% TUDaThesis -- Template für Abschlussarbeiten im CD der TU Darmstadt,
      author=Max Mustermann
    }

    %%%% Für den Titel sind 3 Zeilen vorgesehen.
    %%%% Hat der Titel nur zwei Zeilen, ist ein Zeillenumbruch (\\) einzufügen, um den Titel oben bündig zu positionieren
    %%%% Ist der Titel einzeilig sind zwei Freizeilen (\\ \ \\) anzuhängen.
    %%%% Ludwig Lind 03.02.2021
    \title{Shaping of uniform intensity profiles for single-site addressing of $^{\textbf{85}}\text{Rb}$ atoms}
    \subtitle{Strahlformung von homogenen Intensitätsprofilen für die Einzelplatzadressierung von $^{\textbf{85}}\text{Rb}$-Atomen}
    \author[Daniel Hagn]{Daniel Hagn}%optionales Argument ist die Signatur,
    %\birthplace{Geburtsort}%Geburtsort, bei Dissertationen zwingend notwendig
    \reviewer{Prof. Dr. Gerhard Birkl \and Lukas Sturm}%Gutachter

    %Diese Felder erden untereinander auf der Titelseite platziert.
    %\department ist eine notwendige Angabe, siehe auch dem Abschnitt `Abweichung von den Vorgaben für die Titelseite'
    %\department{Fachbereich Physik} % Das Kürzel wird automatisch ersetzt und als Studienfach gewählt, siehe Liste der Kürzel im Dokument.
    %\institute{Institut für angewandte Physik}
    %\group{Atome-Photonen-Quanten}

    \submissiondate{03.11.2026}
    \examdate{19.11.2026} %\today
    \thesisdate{03.11.2026}

    %    \tuprints{urn=1234,printid=12345}
    %    \dedication{Für alle, die \TeX{} nutzen.}

    \maketitle

    \frontmatter

    %Füge in Abgaben mit twoside=false (z.B. Proposal) trotzdem titleback ein
    \makeatletter
    \if@twoside
    \else%
    \begingroup
    \begin{minipage}[t]{\textwidth}
        \@uppertitleback
    \end{minipage}\par
    \vfill
    \begin{minipage}[b]{\textwidth}
        \@lowertitleback
    \end{minipage}\par
    \@thanks\let\@thanks\@empty
    \endgroup
    \clearpage
    \fi%
    \makeatother


    \tableofcontents
    \printglossary[
    type=abbreviations,
    title=List of Abbreviations,
    nonumberlist]
    
    \mainmatter
  
    \setcounter{page}{1}
    \chapter{Introduction}
    \label{sec:einleitung}
    \input{Content/Introduction}
    
    \chapter{Theoretical Background}
    \label{sec:theory}
    \input{Content/Theory}
    
    \chapter{Test Setup with a Beam Shaper}
    \label{sec:BeamShaperSetup}
    \input{Content/Test Setup}

    \chapter{Multi-tone Simulation for a Top Hat profile}
    \label{sec:constant_simulation}
    \input{Content/Constant_Simulation}
    
    \chapter{Simulation of the beating}
    \label{sec:beating_simulation}
    \input{Content/Beating}
    
    \chapter{New setup}
    \label{sec:newsetup}
    \input{Content/NewSetup}
    
    \chapter{Conclusion and outlook}
    \label{sec:fazit}
    \input{Content/Conclusion}

    \newpage
    
    \appendix
    
    \chapter*{Appendix}
    \addcontentsline{toc}{chapter}{Appendix} % Fügt den Anhang zum Inhaltsverzeichnis hinzu
    \clearpage % Neue Seite vor dem ersten Kapitel im Anhang
    \input{Content/Appendix}

    \chapter*{Friedhof}
    \input{Content/Storage}
    
    \printbibliography
	Internet sources were last accessed on November 3, 2026.
    %\chapter*{Danksagung}
    %Hier kann man Danke sagen.
    %Danke an alle.
    \cleardoublepage
    \affidavit

\end{document}
